% !PRINTED 35
\weberpage{35}
Die Kombination der beiden Transformationen zweiten Grades gibt also die
Multiplikation mit 2.

\begin{center}
§ 10. \gesperrt{Die Transformation dritten Grades.}
\end{center}

Als zweites Beispiel betrachten wir die Transformation dritten Grades.
Die Gleichung

\weq{1}{UV(U-V)(U-\lambda^2V) = T^2\xi\eta(\xi-\eta)(\xi-\varkappa^2\eta)}

fordert, daß jeder der vier Faktoren dritten Grades $U$, $V$, $U-V$,
$U-\lambda^2V$ durch einen der Linearfaktoren $\xi$, $\eta$, $\xi-\eta$,
$\xi-\varkappa^2\eta$ teilbar und daß der Quotient ein Quadrat sei. Wir
setzen also

\weq{2}{U = \xi(a\xi+b\eta)^2, \qquad V = \eta(a'\xi+b'\eta)^2,}

und verlangen noch, daß

\weq{3}{\frac{U-V}{\xi-\eta}, \qquad \frac{U-\lambda^2V}{\xi-\varkappa^2\eta}}

die Quadrate linearer Funktionen werden.

Von den beiden letzten Forderungen folgt die eine aus der anderen, wenn
wir $U$, $V$ so einrichten, daß, von konstanten Faktoren abgesehen, $U$
und $V$ ineinander übergehen durch die Vertauschung von $\xi$, $\eta$ mit
$\varkappa^2\eta$, $\xi$. Durch diese Vertauschung geht aber

\[
\xi(a\xi+b\eta)^2, \qquad \eta(a'\xi+b'\eta)^2
\]

über in

\[
\varkappa^2\eta(b\xi+a\varkappa^2\eta)^2, \qquad \xi(b'\xi+a'\varkappa^2\eta)^2,
\]

und unsere Forderung ist erfüllt, wenn

\weq{4}{\varkappa^2 = \frac{bb'}{aa'}}

ist. Wenn dann

\weq{5}{\lambda^2 = \varkappa^2\left(\frac{ab}{a'b'}\right)^2}

ist, so geht durch diese Vertauschung

\weq{6}{\frac{U-V}{\xi-\eta} \ \text{in} \ \frac{b'^2}{a^2}\,\frac{U-\lambda^2V}{\xi-\varkappa^2\eta}}

über. Nachdem dies festgesetzt, ist nur noch die Bedingung zu erfüllen,
daß die erste der beiden Größen (3) das Quadrat einer linearen Funktion
wird, die wir mit $(\alpha^2\xi-\beta^2\eta)$ bezeichnen, also:

\weq{7}{U-V = (\xi-\eta)(\alpha^2\xi-\beta^2\eta)^2}

oder

\[
U-\xi(\alpha^2\xi-\beta^2\eta)^2 = V-\eta(\alpha^2\xi-\beta^2\eta)^2,
\]
%CONT>>
